% ============================================================
%  示例章节 —— 展示模板提供的各类环境用法与公式规范
%  复制本文件为「第X章.tex」，替换正文即可；用完后在
%  课程笔记模板.tex 中通过 \input{} 引入。
%  数学规范：行间公式一律用 equation 环境并加 \label；
%            向量/矩阵用 \boldsymbol，标量不加粗。
% ============================================================
\section{基本任务}
在实际工程中，构件应满足如下需求：
\begin{myTheorem}{基本任务}
	\begin{enumerate}
	\item 强度要求：在规定载荷作用下的构件不应破坏。
	\item 刚度要求：有足够的抵抗变形的能力。
	\item 稳定性要求：有足够的保持原有平衡形态的能力。
	\end{enumerate}
\end{myTheorem}

\section{基本假设}
为了研究方便，我们引入一些基本假设来简化模型：
\begin{myhypothesis}{基本假设}

\begin{enumerate}
	\item 连续性假设: 认为组成固体的物质不留间隙地充满了固体的体积。
	\item 均匀性假设: 认为固体内处处有相同的力学性能。\\
	\item 各向同性假设: 认为固体的力学性能与方向无关。\\
	\item 小变形假设: 认为变形的尺寸远小于原尺寸。
\end{enumerate}
\end{myhypothesis}

\section{定理与例题}
\begin{theorem}[定理标题]
	这里是定理内容，例如 $a^2 + b^2 = c^2$。
\end{theorem}
\begin{proof}
	这里是证明过程。向量用 $\boldsymbol{a}$，矩阵用 $\boldsymbol{A}$，标量分量用 $a_i$。
\end{proof}

\begin{definition}
	这里是定义内容。
\end{definition}

\begin{example}
	这里是例题内容。
\end{example}
\begin{solution}
	这里是解题过程。
\end{solution}

\begin{remark}
	这里是注释内容。
\end{remark}

\section{公式写法规范}
行间公式一律使用 \texttt{equation} 环境并加 \texttt{\textbackslash label}，禁止裸 \texttt{\$\$} 与手写编号：
\begin{equation}
	\boldsymbol{A}^{\mathrm{T}}\boldsymbol{A}=\boldsymbol{I}
	\label{eq:示例:正交}
\end{equation}
交叉引用用 \texttt{\textbackslash eqref}，例如由式~\eqref{eq:示例:正交} 可知 $\boldsymbol{A}$ 为正交矩阵。向量与矩阵用 $\boldsymbol{}$ 粗体，标量不加粗：
\begin{equation}
	\boldsymbol{v}_{|A}=\boldsymbol{A}\,\boldsymbol{v}_{|B}
	\label{eq:示例:坐标变换}
\end{equation}

\section{公式框}
\begin{myformula}
	$E = mc^2$
\end{myformula}
